\documentclass[11pt]{article}
% preamble.tex --- shared packages, macros, and theorem environments.
% % preamble.tex --- shared packages, macros, and theorem environments.
% % preamble.tex --- shared packages, macros, and theorem environments.
% \input{preamble} after \documentclass in each standalone note.

\usepackage{amsmath, amssymb, amsthm}
\usepackage{mathtools}
\usepackage{geometry}
\usepackage{hyperref}
\usepackage{natbib}
\usepackage{graphicx}
\usepackage{booktabs}
\usepackage{algorithm}
\usepackage{algpseudocode}
\usepackage{xcolor}

\geometry{margin=1in}

\newtheorem{proposition}{Proposition}
\newtheorem{remark}{Remark}

\DeclareMathOperator{\tr}{tr}
\DeclareMathOperator{\diag}{diag}
\newcommand{\R}{\mathbb{R}}
\newcommand{\E}{\mathbb{E}}
\newcommand{\Var}{\mathrm{Var}}
\newcommand{\dd}{\,\mathrm{d}}
\newcommand{\bx}{\mathbf{x}}
\newcommand{\bs}{\mathbf{s}}
\newcommand{\bphi}{\boldsymbol{\phi}}
 after \documentclass in each standalone note.

\usepackage{amsmath, amssymb, amsthm}
\usepackage{mathtools}
\usepackage{geometry}
\usepackage{hyperref}
\usepackage{natbib}
\usepackage{graphicx}
\usepackage{booktabs}
\usepackage{algorithm}
\usepackage{algpseudocode}
\usepackage{xcolor}

\geometry{margin=1in}

\newtheorem{proposition}{Proposition}
\newtheorem{remark}{Remark}

\DeclareMathOperator{\tr}{tr}
\DeclareMathOperator{\diag}{diag}
\newcommand{\R}{\mathbb{R}}
\newcommand{\E}{\mathbb{E}}
\newcommand{\Var}{\mathrm{Var}}
\newcommand{\dd}{\,\mathrm{d}}
\newcommand{\bx}{\mathbf{x}}
\newcommand{\bs}{\mathbf{s}}
\newcommand{\bphi}{\boldsymbol{\phi}}
 after \documentclass in each standalone note.

\usepackage{amsmath, amssymb, amsthm}
\usepackage{mathtools}
\usepackage{geometry}
\usepackage{hyperref}
\usepackage{natbib}
\usepackage{graphicx}
\usepackage{booktabs}
\usepackage{algorithm}
\usepackage{algpseudocode}
\usepackage{xcolor}

\geometry{margin=1in}

\newtheorem{proposition}{Proposition}
\newtheorem{remark}{Remark}

\DeclareMathOperator{\tr}{tr}
\DeclareMathOperator{\diag}{diag}
\newcommand{\R}{\mathbb{R}}
\newcommand{\E}{\mathbb{E}}
\newcommand{\Var}{\mathrm{Var}}
\newcommand{\dd}{\,\mathrm{d}}
\newcommand{\bx}{\mathbf{x}}
\newcommand{\bs}{\mathbf{s}}
\newcommand{\bphi}{\boldsymbol{\phi}}


\title{Universal Function Approximators on $[0,1]^D$:\\
Triangulations and Adaptive Tree Bases for the Mat\'ern SPDE}
\author{Jack Heinzel}
\date{Draft --- \today}

\begin{document}
\maketitle

\begin{abstract}
This note develops the finite-element bases used to discretise the Mat\'ern SPDE
prior of the companion \emph{Foundations} note.  All three are \emph{universal
function approximators} on $[0,1]^D$, but they differ sharply in how they scale
with dimension and how they compose with adaptive refinement.  We treat, in turn:
the \emph{adaptive hat tree basis} (the primary method --- a hierarchical
tensor-product basis of constant, linear, and dyadic hat functions, with a
diagonal stiffness matrix and a recursive-block mass matrix that refine locally);
the \emph{adaptive piecewise-constant tree basis} (a leaf-only indicator basis
whose mass matrix is diagonal and whose operator must be recast in finite-volume
/ GMRF form); and the \emph{Delaunay-triangulation baseline} (Level 1, retained
for context and superseded for the reasons in the lessons appendix).  A common
theme --- the non-uniqueness of the tree label in $D\ge2$, and its canonical
fix --- recurs for both tree bases.  We close with the stationary anisotropic
extension and a lessons-learned appendix.
\end{abstract}

\tableofcontents
\bigskip

% ============================================================
\section{Adaptive Hierarchical Basis on $[0,1]^D$}
\label{sec:hat-tree}
% ============================================================

This is the primary discretisation used throughout this note.  Rather than
meshing the domain with simplices (the Level-1 baseline of
Section~\ref{sec:baseline}, which does not scale to $D\ge 3$; see
Appendix~\ref{sec:lessons}), we represent the field directly in a
hierarchical tree of basis functions adapted to the observed point pattern.
Unlike the leaf-only approach of a rectangular partition, the active basis
includes nodes at \emph{every level} of the tree — the root, its children, and
all their descendants.  This is essential: if only leaf hat-functions were used,
the field would be identically zero at every cell boundary, which prevents
representation of smooth functions.  Including the constant and linear ancestors
remedies this and yields a basis that spans piecewise-linear functions at all
dyadic scales.

\subsection{One-Dimensional Hierarchy}
\label{sec:1d-hierarchy}

Let $\mathcal{D} = [0,1]^D$.  In each dimension $d$ independently, the
basis functions follow a fixed child rule.

\paragraph{Level 0 --- constant.}
\begin{equation}
  f^{(0)}(x) = 1, \quad x \in [0,1].
  \label{eq:f0}
\end{equation}

\paragraph{Level 1 --- linear.}
The unique child of the constant is
\begin{equation}
  f^{(1)}(x) = x, \quad x \in [0,1].
  \label{eq:f1}
\end{equation}

\paragraph{Level 2 --- unit hat.}
The unique child of the linear function is the piecewise-linear hat
\begin{equation}
  f^{(2)}(x) = \hat(x;\,0,1), \qquad
  \hat(x;\,lo,hi) =
  \begin{cases}
    0                  & x \notin [lo,hi], \\
    2(x-lo)/(hi-lo)   & x \in [lo,\tfrac{lo+hi}{2}], \\
    2(hi-x)/(hi-lo)   & x \in [\tfrac{lo+hi}{2},hi].
  \end{cases}
  \label{eq:hat1d}
\end{equation}

\paragraph{Levels $\ge 3$ --- dyadic sub-hats (binary branching).}
From level~2 onward, each hat $\hat(\cdot;\,lo,hi)$ has two children at
the midpoint $m = (lo+hi)/2$:
\begin{equation}
  \hat(\cdot;\,lo,\,m) \quad \text{(left child)}, \qquad
  \hat(\cdot;\,m,\,hi) \quad \text{(right child)}.
  \label{eq:hat-children}
\end{equation}
Every interval is dyadic: $[p/2^n,\,(p+1)/2^n]$ for integers $p,n \ge 0$.

\paragraph{Spanning property.}
The collection $\{1,\, x,\, \hat(\cdot;0,1),\, \hat(\cdot;0,\tfrac12),\,
\hat(\cdot;\tfrac12,1),\ldots\}$ is the \emph{Schauder basis} for $C([0,1])$:
every continuous function on $[0,1]$ can be uniformly approximated by finite
linear combinations.  Including the constant and linear functions is essential;
without them, the span contains only functions vanishing at 0 and 1.

\subsection{$D$-Dimensional Basis and Active Set}
\label{sec:Ddim-basis}

A $D$-dimensional basis function is a product of $D$ one-dimensional
components:
\begin{equation}
  \phi_k(x) = \prod_{d=1}^D f_{s_d^{(k)}}(x_d),
  \label{eq:Ddim-phi}
\end{equation}
where $s_d^{(k)}$ is the 1D state of node $k$ in dimension $d$.
The root has all dimensions at ``const'': $\phi_{\rm root}(x) = 1$.
A \emph{birth move} advances exactly one dimension one step in its local
hierarchy (const $\to$ linear $\to$ unit hat $\to$ left or right sub-hat
$\to \cdots$).

\textbf{The active basis $\{\phi_k\}_{k\in\mathcal{A}}$ consists of all nodes
in the tree --- root, internal nodes, and leaves alike.}
Downward closure holds: if node $k$ is active, so is its parent.

\begin{remark}
If only leaf nodes were used as basis functions, every function in the span
would be zero on every cell boundary (hat functions vanish at the boundary of
their support box).  Including the constant and linear ancestors restores full
approximation power.
\end{remark}

\subsection{Uniqueness: Product Poset versus Birth Tree}
\label{sec:uniqueness}

The labelling $\phi_k$ by a \emph{tree node} $k$ obscures a redundancy that
appears as soon as $D\ge 2$: two distinct nodes of the birth tree can carry the
\emph{same} basis function.  This subsection makes the redundancy precise,
identifies its source (birth moves in distinct dimensions commute), and gives a
canonical labelling that removes it.

\paragraph{The one-dimensional hierarchy as a tree.}
Collect the 1D states of Section~\ref{sec:1d-hierarchy} into a set $\mathcal{H}$,
the nodes of the rooted tree generated by the child rule
\eqref{eq:f0}--\eqref{eq:hat-children}.  Each $\eta\in\mathcal{H}$ carries a
\emph{level}
\begin{equation}
  \ell(\eta) \in \{0,1,2,\dots\},
\end{equation}
the number of advance-steps (births) from the root: $\ell=0$ for the constant
$f^{(0)}$, $\ell=1$ for the linear $f^{(1)}$, $\ell=2$ for the unit hat
$f^{(2)}=\hat(\cdot;0,1)$, and for $\ell\ge 2$ a dyadic \emph{address}
\begin{equation}
  \alpha(\eta)\in\{L,R\}^{\ell-2}
\end{equation}
recording the left/right descent \eqref{eq:hat-children} from the unit hat down
to $\eta$.  The pair $(\ell,\alpha)$ determines $\eta$, hence the 1D function
$f_\eta$.  The advance (successor) relation is \emph{unary} for $\ell<2$ (const
$\to$ linear $\to$ unit hat) and \emph{binary} for $\ell\ge2$ (the two dyadic
children).  Crucially, $\mathcal{H}$ is a tree: the descent from the root to any
$\eta$ is the \emph{unique} path, a word
\begin{equation}
  \rho(\eta)
  = \bigl(\text{lin},\,\text{hat},\,\alpha_1,\dots,\alpha_{\ell-2}\bigr)
  \quad\text{of length } \ell(\eta),
  \label{eq:descent-word}
\end{equation}
whose first $\min(\ell,2)$ steps are forced and whose remaining $\ell-2$ steps
are the address letters.

\paragraph{Basis functions are indexed by a product poset.}
A $D$-dimensional basis function depends only on the \emph{state vector}
\begin{equation}
  s = (s_1,\dots,s_D) \in \mathcal{H}^D,
  \qquad
  \phi_s(x) = \prod_{d=1}^D f_{s_d}(x_d),
  \label{eq:phi-state}
\end{equation}
which refines the node-indexed form \eqref{eq:Ddim-phi}: there $s_d=s_d^{(k)}$
is read off node $k$.  Order $\mathcal{H}^D$ componentwise by the ancestor
relation of $\mathcal{H}$; the root state is $\mathbf{0}=(0,\dots,0)$ with
$\phi_{\mathbf 0}\equiv 1$.

\begin{proposition}[Faithful labelling]
\label{prop:phi-injective}
The map $s\mapsto\phi_s$ from $\mathcal{H}^D$ to functions on $[0,1]^D$ is
injective; indeed the family $\{\phi_s\}_{s\in\mathcal{H}^D}$ is linearly
independent.
\end{proposition}
\begin{proof}
The 1D functions $\{f_\eta\}_{\eta\in\mathcal{H}}$ are linearly independent
(the Faber--Schauder system: constant, linear, and dyadic hats of strictly
increasing resolution).  A tensor product of linearly independent families is
linearly independent, so distinct state vectors give distinct functions.
\end{proof}

Thus the \emph{state vector $s$ is the canonical, non-redundant label} for a
basis function.  The active set is honestly a finite, downward-closed subset
(order ideal) $\mathcal{A}\subset\mathcal{H}^D$: if $s\in\mathcal{A}$ and
$s'\preceq s$ then $s'\in\mathcal{A}$ (Eq.~\eqref{eq:phi-state} with the
spanning requirement of Section~\ref{sec:Ddim-basis}).

\paragraph{The birth tree and its projection.}
The RJMCMC sampler (companion \emph{Sampling Methods} note) does not build
$\mathcal{A}$ directly; it grows a \emph{birth tree} $\mathcal{T}$ whose nodes
are sequences
of birth moves.  A birth move is a pair
\begin{equation}
  b = (d,\,c),\qquad d\in\{1,\dots,D\},
\end{equation}
that advances coordinate $d$ one step in $\mathcal{H}$; the choice $c$ is vacuous
when $\ell(s_d)<2$ and is $c\in\{L,R\}$ when $\ell(s_d)\ge2$.  A node of
$\mathcal{T}$ is a finite birth path $\beta=(b_1,\dots,b_n)$ applied to
$\mathbf 0$, and
\begin{equation}
  \pi:\;\mathcal{T}\longrightarrow\mathcal{H}^D,
  \qquad
  \pi(\beta) = b_n\circ\cdots\circ b_1(\mathbf 0),
\end{equation}
sends a birth path to the state vector it produces.  The node-indexed basis
$\phi_k$ is the composite $k\mapsto\pi(k)\mapsto\phi_{\pi(k)}$.

\begin{proposition}[Commutation of cross-dimensional births]
\label{prop:commute}
If birth moves $b=(d,c)$ and $b'=(d',c')$ act on distinct coordinates
$d\ne d'$, then $b'\circ b = b\circ b'$.
\end{proposition}
\begin{proof}
$b$ alters only the $d$-th component of the state vector and $b'$ only the
$d'$-th; with $d\ne d'$ the two maps act on disjoint coordinates and commute.
\end{proof}

This is the formal content of ``a child of $X\!-\!Y$ equals a child of
$Y\!-\!X$'': advancing dimension $X$ then $Y$ reaches the same state vector as
advancing $Y$ then $X$, hence the same $\phi_s$.  Within a single dimension
there is no such freedom --- the descent \eqref{eq:descent-word} is unique --- so
all redundancy is cross-dimensional interleaving.

\begin{proposition}[The fibre of $\pi$ is a set of shuffles]
\label{prop:fibre}
Fix $s=(s_1,\dots,s_D)\in\mathcal{H}^D$ with levels $\ell_d=\ell(s_d)$.  The
birth paths reaching $s$ are exactly the \emph{shuffles} (order-preserving
interleavings) of the $D$ descent words $\rho(s_1),\dots,\rho(s_D)$ of
\eqref{eq:descent-word}, each step tagged by its dimension.  Hence
\begin{equation}
  \bigl|\pi^{-1}(s)\bigr|
  = \binom{\ell_1+\cdots+\ell_D}{\ell_1,\;\ell_2,\;\dots,\;\ell_D}
  = \frac{(\sum_d \ell_d)!}{\prod_d \ell_d!}.
  \label{eq:fibre-count}
\end{equation}
\end{proposition}
\begin{proof}
A birth path reaching $s$ must, restricted to coordinate $d$, replay the unique
descent $\rho(s_d)$ (Proposition~\ref{prop:phi-injective} and the tree structure
of $\mathcal{H}$); the only freedom is the order in which the $\sum_d\ell_d$
steps from different coordinates are interleaved, and by
Proposition~\ref{prop:commute} every interleaving yields the same $s$.  The
number of interleavings of words of lengths $\ell_1,\dots,\ell_D$ that preserve
each word's internal order is the multinomial coefficient.
\end{proof}

\begin{remark}[Magnitude of the redundancy]
For $D=1$, \eqref{eq:fibre-count} equals $1$: the tree label and the state
vector coincide, and the basis is automatically unique.  For $D\ge2$ it is
$>1$ whenever at least two coordinates have advanced: the user's example
$\ell_X=\ell_Y=1$ gives $\binom{2}{1,1}=2$ (the orders $XY$ and $YX$), and a
state with all $D$ coordinates at the unit-hat level gives $D!$ distinct birth
paths for one basis function.  Counting tree nodes therefore badly
over-counts the basis, and a naive RJMCMC birth proposal visits the same
$\phi_s$ through many paths.
\end{remark}

\paragraph{Canonical labelling (the fix).}
To restore a bijection, impose the fixed total order $1<2<\cdots<D$ on
dimensions and keep, from each fibre, the \emph{unique sorted shuffle} --- the
birth path in which dimension indices are non-decreasing (all dimension-$1$
births, then all dimension-$2$ births, $\dots$).  Equivalently, define the
\emph{canonical parent} of $s\ne\mathbf 0$ by demoting one step in the
highest-indexed advanced coordinate,
\begin{equation}
  \mathrm{par}(s) = \text{demote } s_{d^\star},
  \qquad
  d^\star = \max\{\,d : \ell(s_d)>0\,\}.
  \label{eq:canon-parent}
\end{equation}
Iterating \eqref{eq:canon-parent} drains coordinates from high index to low and
always reaches $\mathbf 0$, so it defines a spanning tree of the Hasse diagram
of the order ideal $\mathcal{A}$; on this tree $\pi$ is a bijection onto
$\mathcal{A}$.

The rule is purely \emph{state-local}, hence usable inside the sampler: a birth
in dimension $d$ is admissible iff every higher coordinate is still at the root,
\begin{equation}
  \text{birth in } d \text{ allowed}
  \iff
  \ell(s_{d'}) = 0 \ \text{for all } d'>d.
  \label{eq:canon-birth-rule}
\end{equation}
With \eqref{eq:canon-birth-rule} the RJMCMC birth/death moves (companion
\emph{Sampling Methods} note) traverse each $\phi_s$ through exactly one path, so
the move counts $N_m,M_m$ and the acceptance ratio no longer double-count
interleaved relabellings of the same basis function.

\subsection{One-Dimensional Inner Products}
\label{sec:1d-inners}

Since $\phi_k$ is a product of 1D functions, all mass and stiffness entries
factor into products of 1D inner products.  Define
\begin{align}
  m(f, g) &= \int_0^1 f(x)\,g(x)\,\dd x, \label{eq:mfg} \\
  s(f, g) &= \int_0^1 f'(x)\,g'(x)\,\dd x. \label{eq:sfg}
\end{align}
Table~\ref{tab:1d-inners} gives closed-form values for all pairings;
$w = hi-lo$ is the hat width.

\begin{table}[h]
\centering
\caption{One-dimensional mass ($m$) and stiffness ($s$) inner products.
${}^\dagger$ by symmetry ($f \leftrightarrow g$).}
\label{tab:1d-inners}
\renewcommand{\arraystretch}{1.35}
\begin{tabular}{ll@{\quad}l}
\toprule
$f$ & $g$ & $m(f,g)$ \\
\midrule
$1$                  & $1$                  & $1$ \\
$1$                  & $x$                  & $\tfrac{1}{2}$ \\
$1$                  & $\hat(\cdot;lo,hi)$  & $\tfrac{w}{2}$ \\
$x$                  & $x$                  & $\tfrac{1}{3}$ \\
$x$                  & $\hat(\cdot;lo,hi)$  & $lo\cdot\tfrac{w}{2} + \tfrac{w^2}{6}$ \\
$\hat(\cdot;lo,hi)$  & $\hat(\cdot;lo,hi)$  & $\tfrac{w}{3}$ \\
$\hat(\cdot;lo,hi)$  & $\hat(\cdot;lo,m)$\,${}^\dagger$
                                             & $\tfrac{w}{8}$ \\
$\hat(\cdot;lo,m)$   & $\hat(\cdot;m,hi)$   & $0$ \\
\bottomrule
\end{tabular}
\qquad
\begin{tabular}{ll@{\quad}l}
\toprule
$f$ & $g$ & $s(f,g)$ \\
\midrule
$1$                  & $\cdot$ (any)         & $0$ \\
$x$                  & $x$                   & $1$ \\
$x$                  & $\hat(\cdot;lo,hi)$   & $0$ \\
$\hat(\cdot;lo,hi)$  & $\hat(\cdot;lo,hi)$   & $\tfrac{4}{w}$ \\
$\hat(\cdot;lo,hi)$  & $\hat(\cdot;lo,m)$\,${}^\dagger$
                                              & $0$ \\
$\hat(\cdot;lo,m)$   & $\hat(\cdot;m,hi)$    & $0$ \\
\bottomrule
\end{tabular}
\end{table}

\noindent
Key observations:
\begin{itemize}
  \item $s(1,\cdot) = 0$: the constant has zero derivative, contributing
        nothing to the stiffness matrix.
  \item $s(x, \hat) = 0$: by the fundamental theorem,
        $\int_{lo}^{hi} 1 \cdot \hat'\,\dd x = \hat(hi) - \hat(lo) = 0$,
        so the linear function is stiffness-orthogonal to every hat.
  \item $s(\hat(lo,hi),\,\hat(lo,m)) = 0$: on $[lo,m]$ (the support of
        $\hat_L$), the parent derivative $\hat'(x;lo,hi) = 2/w$ is constant,
        so $\int_{lo}^m (2/w)\,\hat'_L\,\dd x = (2/w)[\hat_L(m)-\hat_L(lo)]
        = 0$.
  \item $s(\hat_L,\hat_R) = 0$: disjoint supports.
\end{itemize}
Consequently, \emph{all off-diagonal stiffness inner products vanish}: the
stiffness matrix is diagonal.

\subsection{Mass Matrix}
\label{sec:hat-mass}

\begin{proposition}[Tensor-product mass matrix]
\label{prop:mass-new}
For any two active basis functions $\phi_i$, $\phi_j$:
\begin{equation}
  C_{ij} = \int_{[0,1]^D} \phi_i(x)\,\phi_j(x)\,\dd x
          = \prod_{d=1}^D m\!\bigl(f_{s_d^{(i)}},\, f_{s_d^{(j)}}\bigr).
  \label{eq:hat-mass-new}
\end{equation}
\end{proposition}
\begin{proof}
$\phi_i\phi_j = \prod_d f_{s_d^{(i)}}(x_d)f_{s_d^{(j)}}(x_d)$.
Fubini gives the product of 1D integrals.
\end{proof}

\paragraph{Sparsity.}
$C_{ij} = 0$ when $m(f_{s_d^{(i)}},f_{s_d^{(j)}})=0$ for some $d$.
From the table, this occurs when the dimension-$d$ states are sibling hats
(zero support overlap).  In particular, the two children of any hat split
have zero mass coupling.

\paragraph{Block structure.}
When node $k$ is split in dimension $d_0$, its children $k_L$ and $k_R$
carry sibling hat states in that dimension, so $C_{ij} = 0$ for all
$i\in\mathcal{A}_{k_L}$, $j\in\mathcal{A}_{k_R}$:
\begin{equation}
  C_{\mathcal{A}_{k_L},\,\mathcal{A}_{k_R}} = 0.
  \label{eq:C-block-zero}
\end{equation}
Ancestors of $k$ (not yet split in $d_0$) couple to both subtrees and lie
outside this block.  The mass matrix therefore has a \emph{recursive
block-triangular} structure with blocks indexed by the subtrees of $\mathcal{T}$.

\subsection{Stiffness Matrix}
\label{sec:hat-stiffness}

\begin{proposition}[Diagonal stiffness matrix]
\label{prop:stiffness-new}
$G$ is diagonal:
\begin{equation}
  G_{ij} = \delta_{ij}
    \sum_{d=1}^D s\!\bigl(f_{s_d^{(i)}},\, f_{s_d^{(i)}}\bigr)
    \prod_{d' \neq d} m\!\bigl(f_{s_{d'}^{(i)}},\, f_{s_{d'}^{(i)}}\bigr).
  \label{eq:hat-stiffness-new}
\end{equation}
\end{proposition}
\begin{proof}
The general formula is $G_{ij} = \sum_d s(f_{s_d^{(i)}},f_{s_d^{(j)}})
\prod_{d'\neq d}m(f_{s_{d'}^{(i)}},f_{s_{d'}^{(j)}})$.  Every off-diagonal
$s$-factor vanishes (as shown in Section~\ref{sec:1d-inners}), so only the
$i=j$ term survives.
\end{proof}

\noindent
The diagonal entry $G_{kk}$ is the sum of $D$ terms, each equal to the
stiffness self-product in one dimension times the mass self-product in
all other dimensions.  For the root ($f_{s_d} = 1$ for all $d$), $G_{kk}=0$.
For a node with all dimensions at ``const'' except dimension $d_0$ at
$\hat(\cdot;lo,hi)$ (width $w$):
\begin{equation}
  G_{kk} = \frac{4}{w} \cdot \prod_{d \neq d_0} m(f_{s_d},f_{s_d}).
  \label{eq:Gkk-example}
\end{equation}

\subsection{SPDE Operator and Prior}
\label{sec:hat-spde}

The FEM operator $K = \kappa^2 C + G$ combines the (sparse) mass matrix with
the diagonal stiffness.  The $\alpha=2$ SPDE precision is
\begin{equation}
  Q = K C^{-1} K.
  \label{eq:hat-Q-new}
\end{equation}
$C$ is sparse and positive definite, so $C^{-1}$ can be applied via sparse
Cholesky; $K^{-1}$ (needed for HMC velocity refresh and CG solves) inherits
the same sparsity.  The prior is $r \sim \mathcal{N}(0,\,\sigma^2 Q^{-1})$
over all $|\mathcal{A}|$ active coefficients simultaneously.

\paragraph{Likelihood.}
For each observation $s_i$, the log-intensity sums over all active basis
functions:
\begin{equation}
  \log\lambda(s_i) = \mu + \bphi(s_i)^T r,
  \quad
  \bphi(s_i)_k = \prod_{d=1}^D f_{s_d^{(k)}}(s_{i,d}).
  \label{eq:hat-loglik-new}
\end{equation}
The constant contributes 1, the linear terms contribute the observation
coordinates, and at most one hat per dimension is nonzero at $s_i$.
The Poisson integral is approximated by a fixed quadrature rule.
% ============================================================
\section{Adaptive Piecewise-Constant Tree Basis}
\label{sec:const-tree}
% ============================================================

The hat tree basis of Section~\ref{sec:hat-tree} pays a price for its smoothness:
the active set must include \emph{every} ancestor (root, linears, coarse hats),
and those ancestors overlap broadly --- the root constant overlaps every cell,
each linear overlaps its whole slab --- so the mass matrix acquires dense rows as
$D$ grows.  This section explores the opposite end of the regularity spectrum:
keep the same adaptive dyadic refinement of $[0,1]^D$, but take the basis
functions to be \emph{indicators of the leaf cells}.  The active set is then the
\emph{leaves only}, and --- as we show --- this still yields a universal function
approximator while making the mass matrix exactly diagonal.  The cost is that the
SPDE operator can no longer be discretised by a conforming Galerkin stiffness
matrix; it must be recast in finite-volume / GMRF form, which we develop and
evaluate below.

\subsection{Partition and Basis}
\label{sec:const-partition}

We refine $[0,1]^D$ with the same binary kd-tree as before: a cell is split at
the midpoint of one axis into two children.  At any stage the \emph{leaves}
$\{C_i\}_{i\in\mathcal{L}}$ form a partition of the domain into disjoint
axis-aligned dyadic boxes,
\begin{equation}
  C_i = \prod_{d=1}^D I_d^{(i)},
  \qquad
  \bigcup_{i\in\mathcal{L}} C_i = [0,1]^D,
  \quad
  C_i \cap C_j = \emptyset \ (i\ne j),
  \label{eq:const-partition}
\end{equation}
where each $I_d^{(i)}=[p/2^{n},(p+1)/2^n]$ is a dyadic subinterval.  The basis is
the set of leaf indicators
\begin{equation}
  \phi_i(x) = \mathbf{1}[x\in C_i],
  \qquad i\in\mathcal{L}.
  \label{eq:const-basis}
\end{equation}
As with the hat tree, a leaf carries the product structure
$C_i=\prod_d I_d^{(i)}$: it is indexed by a $D$-tuple of 1D dyadic addresses, the
number of splits in dimension $d$ being the depth $n_d^{(i)}=\log_2(1/|I_d^{(i)}|)$.

\subsection{Universal Approximation from Leaves Alone}
\label{sec:const-ufa}

\begin{proposition}[Density of refining indicators]
\label{prop:const-ufa}
Let $\delta(\mathcal{L})=\max_{i\in\mathcal{L}}\operatorname{diam}(C_i)$ be the
mesh size.  As $\delta(\mathcal{L})\to0$, the span of $\{\phi_i\}_{i\in\mathcal{L}}$
is dense in $L^2([0,1]^D)$, and for every $g\in C([0,1]^D)$ the piecewise-constant
interpolant $\sum_i g(\bar x_i)\phi_i$ (with $\bar x_i$ the cell centroid)
converges uniformly to $g$.
\end{proposition}
\begin{proof}
The leaves partition the domain, so $\sum_i\phi_i\equiv1$ (partition of unity).
For $g\in C([0,1]^D)$, uniform continuity gives
$\sup_x|g(x)-\sum_i g(\bar x_i)\phi_i(x)|\le\omega_g(\delta(\mathcal{L}))\to0$,
where $\omega_g$ is the modulus of continuity; uniform approximation of
continuous functions implies $L^2$-density of the span by the density of
$C([0,1]^D)$ in $L^2$.
\end{proof}

\begin{remark}[Why leaves suffice here but not for hats]
This is the structural advantage over the hat tree.  A \emph{leaf hat} vanishes
on the boundary of its support box, so the span of leaf hats alone contains only
functions vanishing on every dyadic grid line --- not dense --- which is exactly
why Section~\ref{sec:Ddim-basis} must include all ancestors.  Indicators instead
\emph{tile} the domain and form a partition of unity, so no ancestors are needed:
the active set is the leaf set, and there is no broadly-overlapping root or
linear term.
\end{remark}

\subsection{Mass Matrix is Diagonal}
\label{sec:const-mass}

\begin{proposition}[Diagonal mass matrix]
\label{prop:const-mass}
For the leaf indicators \eqref{eq:const-basis},
\begin{equation}
  C_{ij} = \int_{[0,1]^D}\phi_i\phi_j\dd x
         = \delta_{ij}\,|C_i|,
  \qquad
  |C_i| = \prod_{d=1}^D |I_d^{(i)}|.
  \label{eq:const-mass}
\end{equation}
\end{proposition}
\begin{proof}
$\phi_i\phi_j=\mathbf{1}[x\in C_i\cap C_j]$, which is $\mathbf{1}[x\in C_i]$ for
$i=j$ and $0$ otherwise because leaves are disjoint \eqref{eq:const-partition}.
The integral of $\mathbf{1}[x\in C_i]$ is the cell volume.
\end{proof}

The mass matrix is therefore diagonal in \emph{every} dimension --- a strict
improvement on the hat tree, whose ancestor functions couple to whole subtrees.
The lumped and consistent mass matrices coincide ($\tilde C=C$), so the
quadrature weights $c_i=|C_i|$ of the likelihood integral
(\emph{Foundations} note, integral-approximation section) are exact cell volumes,
and no buffer-node lumping approximation is incurred.

\subsection{Stiffness: An Obstruction and Its Finite-Volume Resolution}
\label{sec:const-stiffness}

Here the piecewise-constant basis runs into a genuine obstruction, which we state
plainly before resolving it.

\begin{proposition}[The conforming stiffness matrix does not exist]
\label{prop:const-stiff-illdef}
Indicator functions are not in $H^1([0,1]^D)$: the distributional gradient of
$\phi_i=\mathbf 1[x\in C_i]$ is a (vector-valued) surface measure supported on
$\partial C_i$, not an $L^2$ function.  Consequently the Galerkin stiffness
integral $G_{ij}=\int\nabla\phi_i\cdot\nabla\phi_j\dd x$ is undefined (formally
divergent), and the weak form of the \emph{Foundations} note --- which integrates
by parts and requires $\phi_i\in H^1$ --- does not apply.
\end{proposition}

In other words the claim ``the stiffness matrix stays sparse'' cannot be taken at
face value: there is no conforming stiffness matrix to speak of.  What survives
is a \emph{discrete} second-order operator built from cell adjacencies, in the
spirit of finite-volume schemes and Gaussian Markov random fields (GMRFs).

\paragraph{Adjacency and the finite-volume operator.}
Call two leaves $i,j$ \emph{adjacent}, $i\sim j$, if their closures share a
$(D{-}1)$-dimensional face $F_{ij}$ of positive area.  Discretising
$-\nabla\cdot\nabla$ by a two-point flux across each shared face gives the
sparse, symmetric operator
\begin{equation}
  G^{\mathrm{FV}}_{ij} =
  \begin{cases}
    -\,w_{ij}, & i\sim j,\quad
      w_{ij}=\dfrac{|F_{ij}|}{\lVert \bar x_i-\bar x_j\rVert},\\[1.2ex]
    \ \sum_{k\sim i} w_{ik}, & i=j,\\[0.6ex]
    \ 0, & \text{otherwise},
  \end{cases}
  \label{eq:const-fv}
\end{equation}
with $\bar x_i$ the cell centroid and $|F_{ij}|$ the shared-face area.  This is a
weighted graph Laplacian on the cell-adjacency graph; the precision
\begin{equation}
  Q = \kappa^2 C + G^{\mathrm{FV}},
  \qquad C=\diag(|C_i|),
  \label{eq:const-Q}
\end{equation}
is exactly a Besag-type GMRF (a conditional-autoregression), and is sparse and
positive definite for $\kappa>0$.  It coincides with the standard finite-volume
discretisation of $(\kappa^2-\Delta)$ on the dual mesh.

\begin{proposition}[Sparsity is retained in all dimensions]
\label{prop:const-sparse}
With $C$ diagonal \eqref{eq:const-mass} and $G^{\mathrm{FV}}$ coupling only
face-adjacent leaves \eqref{eq:const-fv}, the number of nonzeros of $Q$ is
$|\mathcal{L}| + 2\,\#\{\text{internal faces}\}$, independent of any global
``ancestor'' coupling.  In particular $Q$ has no dense rows, in contrast to the
hat tree, whose root and linear functions couple to $\Theta(|\mathcal{A}|)$ basis
functions each.
\end{proposition}

\paragraph{Evaluation of the claim.}
The original conjecture --- that the mass and stiffness matrices ``do not become
dense as $D$ increases'' --- is therefore \emph{partly true and worth a precise
verdict}:
\begin{itemize}
  \item \textbf{Mass matrix: true, and strictly better than the hat tree.}
        It is diagonal in every dimension (Proposition~\ref{prop:const-mass}).
  \item \textbf{Stiffness matrix: true only after reinterpretation.}  A
        conforming $H^1$ stiffness matrix does \emph{not exist}
        (Proposition~\ref{prop:const-stiff-illdef}); the right object is the
        finite-volume / GMRF operator \eqref{eq:const-fv}, which \emph{is} sparse
        and free of dense ancestor rows (Proposition~\ref{prop:const-sparse}).
\end{itemize}
Three caveats temper the win:
\begin{enumerate}
  \item \textbf{Adaptive adjacency is irregular.}  Under non-uniform refinement a
        coarse leaf can be face-adjacent to many small leaves, so the per-row
        degree of $G^{\mathrm{FV}}$ is not uniformly $2D$; the matrix is sparse
        but its bandwidth depends on the refinement balance.  Finding the
        face-neighbours of a kd-tree leaf is also more work than the hat tree's
        $O(D)$ tensor update.
  \item \textbf{Rougher prior.}  Piecewise-constant fields are discontinuous, so
        the basis represents an intrinsically rough field; the operator
        \eqref{eq:const-Q} is a consistent discretisation only of the roughest
        Mat\'ern regime, and the smoothness parameter $\nu$ loses its
        finite-element interpretation.  (One recovers smoother fields only by
        refining, or by switching to higher-order discontinuous-Galerkin
        elements with interior-penalty face terms, which reintroduce face
        coupling of the same sparsity class.)
  \item \textbf{Same GMRF, different CG constant.}  Because \eqref{eq:const-Q} is
        a Besag CAR precision, the dimension-robust samplers of the
        \emph{Sampling Methods} note apply verbatim (the CG conditioning argument
        there is stated for exactly this $Q=\kappa^2 I + L_{\mathrm{graph}}$
        form).
\end{enumerate}

\subsection{Non-Uniqueness, and the Same Canonical Fix}
\label{sec:const-uniqueness}

The piecewise-constant basis inherits the non-uniqueness of
Section~\ref{sec:uniqueness}.  A leaf box $C_i=\prod_d I_d^{(i)}$ is reached by
cutting dimension $d$ a total of $n_d^{(i)}$ times, and --- by the same
commutation argument (Proposition~\ref{prop:commute}) --- cuts in distinct
dimensions commute: splitting a cell along $X$ and then its children along $Y$
produces the \emph{same} $2\times2$ partition as splitting along $Y$ then $X$.
Hence many distinct kd-trees realise the same leaf partition, and therefore the
same basis; the number of split orders reaching a given leaf is again the
multinomial $\binom{\sum_d n_d}{n_1,\dots,n_D}$ of Eq.~\eqref{eq:fibre-count}.

\begin{remark}[What is redundant here]
For the hat tree the redundancy is over birth \emph{paths} to a single
basis function $\phi_s$; for the constant tree it is over kd-\emph{trees}
generating the same leaf \emph{partition} $\{C_i\}$.  In both cases the honest
index is geometric (the state vector / the leaf box), not the tree.
\end{remark}

The same canonical labelling of Section~\ref{sec:uniqueness} applies, but with
one caveat that does \emph{not} arise for the hat tree.  Imposing the total order
$1<2<\cdots<D$ and keeping only the canonical (sorted) split history --- the
state-local rule that, within a cell's split lineage, a cut in dimension $d$ is
admissible only if no higher-indexed dimension has been cut
(Eq.~\eqref{eq:canon-birth-rule}, equivalently the canonical parent
\eqref{eq:canon-parent}) --- does make the tree-to-partition map injective.  The
subtlety is that, unlike the hat tree where the sorted shuffle reaches
\emph{every} state vector, this restriction also \emph{shrinks the family of
reachable partitions}: once a branch has cut a higher-indexed axis it can never
re-split a lower-indexed one within that subtree, so partitions that genuinely
require interleaved cuts down a single branch become unreachable.  The canonical
rule is therefore a bijection onto the \emph{canonically reachable} partitions, a
strict subfamily, rather than a faithful relabelling of all dyadic-box
partitions.

Two consequences for the RJMCMC sampler follow.  If one is content to explore
only the canonically reachable subfamily, applying
Eq.~\eqref{eq:canon-birth-rule} as a state-local birth restriction removes the
cross-dimensional double-counting at no extra bookkeeping cost.  To retain the
\emph{full} partition family without double-counting, the canonicalisation must
instead be done at the level of the partition itself --- e.g.\ comparing a
proposed split's resulting leaf set against a canonical (say, guillotine-ordered)
representative --- rather than by the per-lineage sorted-cut rule.  In the
implementation the sorted-cut rule is exposed as an opt-in helper for this
reason.

% ============================================================
\section{Baseline (Level 1): Adaptive Triangulation FEM}
\label{sec:baseline}
\label{sec:adaptive-tri}
% ============================================================

This section documents the original Level-1 pipeline: an adaptive kd-tree mesh,
a Delaunay triangulation of its nodes, and piecewise-linear hat functions on the
resulting simplices.  It predates the adaptive tree hat basis of
Section~\ref{sec:hat-tree} and has since been superseded by it, for the reasons
recorded in Appendix~\ref{sec:lessons}: the simplex count grows
super-exponentially with dimension and the triangulation is not static under
refinement.  We retain the construction here because it is the concrete instance
behind the general FEM framework of the companion \emph{Foundations} note, and
because the likelihood, hyperparameter, anisotropic, and RJMCMC machinery were
all first developed on it.

\subsection{Recursive Equal-Count kd-Tree}

We build a mesh that is automatically denser where the observed point pattern is
dense.  Starting from a root cell covering the domain, we recursively split
each cell along its longest axis at the data median until every leaf cell
contains at most $n_{\max}$ points (the refinement threshold).  This is a
\emph{balanced} kd-tree partition: each split divides the data approximately
evenly, so the resulting leaf sizes are $O(n/n_{\rm leaves})$.

The split value is clamped away from the cell boundary by a relative tolerance
$\epsilon = 10^{-8}$ of the cell width, ensuring no degenerate child cell.

\subsection{Mesh Nodes and Delaunay Triangulation}

Each leaf cell contributes one node to the mesh: the centroid of its data
points if the cell is non-empty, or the geometric cell midpoint otherwise.
The $2^d$ corners of the bounding box are added as fixed boundary nodes.

The Delaunay triangulation of all nodes defines the mesh $\mathcal{T}$.
It is rebuilt from scratch after each refinement or coarsening step.

\subsection{Buffer Region}

To ensure the triangulation is regular near the domain boundary, the bounding
box is taken slightly larger than the observation domain---typically
$[-0.1, 1.1]^2$ when observations lie in $[0,1]^2$.  This introduces
\emph{buffer nodes} in the extended region.  We account for them explicitly in
the likelihood (companion \emph{Foundations} note).

\subsection{Basis Functions and Local Geometry}

On the triangulation $\mathcal{T}$ with $m$ nodes and $T$ simplices, we use
piecewise linear \emph{hat} basis functions $\{\phi_i\}_{i=1}^m$.  The hat
function $\phi_j$ equals 1 at node $j$, 0 at every other node, and varies
linearly in between.  Within any single simplex it is a flat tilted plane, so
its gradient is \emph{constant} on each simplex and can be computed from
geometry alone.

\paragraph{Reference simplex and the map $B_\tau$.}
Every computation is carried out on a fixed \emph{reference simplex}
$\hat\tau$ with vertices at the origin $\hat{v}_0=\mathbf{0}$ and the
standard basis vectors $\hat{v}_k = e_k$ for $k=1,\ldots,d$.  On this
template the hat functions are trivially
\[
  \hat\phi_0(\hat s) = 1 - \hat s_1 - \cdots - \hat s_d, \qquad
  \hat\phi_k(\hat s) = \hat s_k \;\;(k\ge 1),
\]
with reference gradients $\nabla_{\hat s}\hat\phi_j = \hat{e}_j$, where
$\hat{e}_0 = -\mathbf{1}$ and $\hat{e}_k$ is the $k$-th standard basis
vector.

The matrix
\begin{equation}
  B_\tau = [v_1 - v_0, \;\ldots\;, v_d - v_0] \in \R^{d\times d}
\end{equation}
maps the reference simplex to the physical simplex $\tau$ via
$s = v_0 + B_\tau\hat s$.  Its columns are the edge vectors of $\tau$
emanating from $v_0$, so $B_\tau$ encodes the shape and orientation of the
triangle.  The volume follows immediately: $|\tau| = |\det B_\tau|/d!$.

\paragraph{Why $B_\tau^{-\top}$ for gradients.}
The gradient $\nabla_s\phi_j^{(\tau)}$ is the vector $g$ satisfying
$g\cdot\delta s = \delta\phi$ for every physical displacement $\delta s$.
Writing $\delta s = B_\tau\,\delta\hat s$ and using the reference gradient
$\nabla_{\hat s}\hat\phi_j = \hat{e}_j$:
\[
  g \cdot B_\tau\,\delta\hat s
  = \hat{e}_j \cdot \delta\hat s
  \quad\text{for all }\delta\hat s,
\]
which gives $B_\tau^{\top}g = \hat{e}_j$, and therefore
\begin{equation}
  \nabla_s \phi_j^{(\tau)} = B_\tau^{-\top}\hat{e}_j, \quad j = 0, \ldots, d.
  \label{eq:grad-phi}
\end{equation}
Intuitively, $B_\tau^{-1}$ converts a physical direction into reference
coordinates, and the transpose accounts for the fact that gradients scale
inversely with distances (a squeezed triangle has steeper slopes).  Because
$B_\tau$ is constant on $\tau$, so is $\nabla_s\phi_j^{(\tau)}$: the hat
function is a flat ramp, and \eqref{eq:grad-phi} gives its slope.

\subsection{Mass Matrix}
\label{sec:mass}

The mass matrix $C \in \R^{m\times m}$ has entries
\begin{equation}
  C_{ij} = \int_\mathcal{D} \phi_i(s)\phi_j(s) \dd s
         = \sum_{\tau \ni i,j} |\tau| \frac{1 + \delta_{ij}}{(d+1)(d+2)}.
  \label{eq:mass}
\end{equation}
$C_{ij}$ measures the \emph{overlap} between hat functions $\phi_i$ and
$\phi_j$: it is nonzero only when nodes $i$ and $j$ share at least one
simplex.  The diagonal entry $C_{ii}$ is proportional to the total area of
all simplices containing node $i$.  The matrix is symmetric positive definite
with sparsity matching the mesh topology.

The \emph{row sum} $\tilde{c}_j = (C\mathbf{1})_j = \sum_k C_{jk}$ is the
total area ``owned'' by node $j$: it equals $\sum_{\tau\ni j}|\tau|/(d+1)$,
the average share of each incident simplex.  Geometrically, $\tilde{c}_j$
approximates the area of the Voronoi cell\footnote{It's not really a Voronoi cell, but it's close-ish to one depending on the triangulation.} around node $j$.  The diagonal
matrix $\tilde{C} = \diag(C\mathbf{1})$ is the \emph{lumped} mass matrix,
used in the companion \emph{Foundations} note to approximate the Poisson integral
as a weighted sum $\sum_j \tilde{c}_j \lambda_j$.

In the SPDE operator $(\kappa^2 - \Delta)$, the term $\kappa^2$ acts like a
scalar multiple of the identity on functions.  After Galerkin projection onto
the hat-function basis, this becomes $\kappa^2 C$: the mass matrix is the
FEM discretisation of the identity operator, weighted by local area.

\subsection{Stiffness Matrix}
\label{sec:stiffness}

The stiffness matrix $G \in \R^{m\times m}$ has entries
\begin{equation}
  G_{ij} = \int_\mathcal{D} \nabla\phi_i(s) \cdot \nabla\phi_j(s) \dd s
         = \sum_{\tau \ni i,j} |\tau| \, \nabla\phi_i^{(\tau)} \cdot \nabla\phi_j^{(\tau)}.
  \label{eq:stiff}
\end{equation}
$G_{ij}$ measures how much nodes $i$ and $j$ \emph{resist spatial variation}
together: it is the inner product of the slopes of the two hat functions.
By integration by parts, $\mathbf{x}^T G\, \mathbf{x} =
\int|\nabla x(s)|^2\dd s$, so the quadratic form in $G$ is exactly the total
squared-gradient energy of the piecewise-linear field
$x(s)=\sum_j x_j\phi_j(s)$.  Smooth fields (small gradients everywhere) have
low $G$-energy; rough or wiggly fields have high $G$-energy.

Equivalently, $G$ discretises the negative Laplacian $-\Delta$: it can be
thought of as a network of \emph{springs} connecting adjacent nodes, where
$G_{ij}$ is (minus) the spring constant for the edge between $i$ and $j$.
For a Delaunay mesh in 2D, $G_{ij} \le 0$ for all neighbours
($i\neq j$), and the action $(G\mathbf{x})_i = -\sum_j G_{ij}(x_j - x_i)$
sums the spring forces pulling node $i$ toward its neighbours' values.

Together, $K = \kappa^2 C + G$ discretises $(\kappa^2 - \Delta)$: the mass
term $\kappa^2 C$ penalises large \emph{absolute} values of the field, and
the stiffness term $G$ penalises spatial \emph{roughness}.  Large $\kappa$
makes the mass term dominate, producing a field with short correlation length
(rapidly decorrelated); small $\kappa$ lets the stiffness term dominate,
producing a smoother field correlated over longer distances.

For the generalised, spatially varying-diffusion case
$(\kappa(s)^2-\nabla\cdot H(s)\nabla)^{\alpha/2}x=\mathcal{W}$ (companion
\emph{Foundations} note), the integrand becomes
$\nabla\phi_i^T H(\bar\tau) \nabla\phi_j$ where $H$ is evaluated at the
simplex centroid $\bar\tau$: the anisotropic diffusion tensor $H$ stretches
the spring network directionally, making correlations longer in some
directions than others.
% ============================================================
\section{Stationary Anisotropic Extension}
\label{sec:anisotropic}
% ============================================================

This extension was developed on the triangulation baseline
(Section~\ref{sec:baseline}), where per-simplex gradients make the
diffusion-tensor assembly natural; the same construction carries over to the
adaptive tree basis through its diagonal stiffness matrix.  We keep the simplex
presentation here because that is where it was first worked out.

\subsection{Anisotropy Matrix Parameterisation}

The isotropic model (companion \emph{Foundations} note) treats all directions
equally through a scalar range parameter $\kappa$.  A stationary
\emph{anisotropic} generalisation retains a constant diffusion tensor $H$ in the
generalised SPDE $(\kappa^2-\nabla\cdot H\nabla)^{\alpha/2}x=\mathcal{W}$
(companion \emph{Foundations} note).  In $\R^2$ we parameterise
\begin{equation}
  H = R(\theta)\,\diag(\kappa_x^2,\,\kappa_y^2)\,R(\theta)^T,
  \label{eq:H-param}
\end{equation}
where $R(\theta)$ is the $2\times 2$ rotation matrix and $\kappa_x, \kappa_y > 0$
are effective range parameters along the principal axes of the field.
The elements of $H$ are
\begin{align}
  h_{11} &= \kappa_x^2\cos^2\!\theta + \kappa_y^2\sin^2\!\theta, \nonumber\\
  h_{12} &= (\kappa_x^2 - \kappa_y^2)\cos\theta\sin\theta,           \label{eq:H-components} \\
  h_{22} &= \kappa_x^2\sin^2\!\theta + \kappa_y^2\cos^2\!\theta.    \nonumber
\end{align}
The three parameters $(\log\kappa_x,\, \log\kappa_y,\, \theta)$ replace the
single scalar $\log\kappa$ from the isotropic model; the marginal variance
$\sigma$ remains a scalar.  The isotropic case is recovered by setting
$\kappa_x = \kappa_y$ (any $\theta$), yielding $H = \kappa_x^2 I$.

\subsection{FEM Basis Decomposition}
\label{sec:aniso-fem}

The anisotropic stiffness matrix
\[
G_{H,ij} = \int_\mathcal{D} \nabla\phi_i^T H\,\nabla\phi_j\,\dd s
\]
decomposes linearly over three basis matrices:
\begin{equation}
  G_H = h_{11}\,G_{xx} + h_{12}\,G_{xy} + h_{22}\,G_{yy},
  \label{eq:GH-decomp}
\end{equation}
with
\begin{align}
  G_{xx,ij} &= \sum_\tau |\tau|\,
    \partial_x\phi_i^{(\tau)}\,\partial_x\phi_j^{(\tau)}, \\
  G_{yy,ij} &= \sum_\tau |\tau|\,
    \partial_y\phi_i^{(\tau)}\,\partial_y\phi_j^{(\tau)}, \\
  G_{xy,ij} &= \sum_\tau |\tau|\!\left(
    \partial_x\phi_i^{(\tau)}\,\partial_y\phi_j^{(\tau)}
    + \partial_y\phi_i^{(\tau)}\,\partial_x\phi_j^{(\tau)}\right).
\end{align}
All three matrices are assembled once from the fixed mesh prior to sampling and
stored in COO sparse format.  The isotropic $G = G_{xx} + G_{yy}$ is the
special case $H = I$.

At each NUTS step, forming $G_H$ costs three sparse scalar-matrix additions, and
the SPDE operator
\begin{equation}
  K(H) = C + G_H(H)
  \label{eq:K-aniso}
\end{equation}
has the same sparsity pattern as the isotropic $K(\kappa)$.  Matrix--vector
products $K(H)\bx$ therefore cost $O(\mathrm{nnz}) \approx O(m)$ via COO
scatter-add, unchanged from the isotropic case.

\subsection{Log-Determinant Challenge}
\label{sec:logdet-aniso}

The companion \emph{Foundations} note showed that for the isotropic operator
$K(\kappa) = \kappa^2 C + G$, the log-determinant reduces to $O(m)$ work per step
once the generalised eigenpairs $G v_i = \mu_i C v_i$ have been precomputed:
\[
  \log|K(\kappa)| = \log|C| + \textstyle\sum_i \log(\kappa^2 + \mu_i).
\]
This is possible because $K$ is a \emph{scalar-affine} function of the single
parameter $\kappa^2$: the eigenvectors $\{v_i\}$ are fixed and only the
eigenvalues shift.

For the anisotropic operator \eqref{eq:K-aniso}, the matrix
$K(H) = C + h_{11}G_{xx} + h_{12}G_{xy} + h_{22}G_{yy}$
is a linear combination of four sparse matrices whose coefficients
$(h_{11}, h_{12}, h_{22})$ all change at every NUTS step.  No fixed eigenbasis
can diagonalise $K(H)$ for all $H$ simultaneously: a generalised eigenproblem
$K(H) v = \mu C v$ would need to be resolved whenever $(h_{11}, h_{12}, h_{22})$
changes.

Consequently, evaluating $\log|K(H)|$ requires a dense factorisation at each
leapfrog step,
\begin{equation}
  \log|K(H)| = \texttt{slogdet}\bigl(K(H)\bigr),
  \quad\text{cost } O(m^3) \text{ per step.}
  \label{eq:slogdet-aniso}
\end{equation}
The gradient required by NUTS,
\begin{equation}
  \frac{\partial\log|K|}{\partial\theta_k}
  = \tr\!\left(K^{-1}\frac{\partial K}{\partial\theta_k}\right),
  \label{eq:grad-logdet-aniso}
\end{equation}
requires an implicit $K^{-1}$ solve for each of the three anisotropy
parameters $(\log\kappa_x, \log\kappa_y, \theta)$, which JAX's automatic
differentiation of \texttt{jnp.linalg.slogdet} provides via the implicit
function theorem, again at $O(m^3)$ cost.

For meshes of size $m \lesssim 500$ the per-step $O(m^3)$ cost is acceptable
($\approx\!1$\,ms on a modern CPU), but it precludes scaling to
$m \gtrsim 2000$.  Finding an efficient $O(m)$--$O(m\log m)$ approximation
to $\log|K(H)|$ for parametric matrix pencils is an \textbf{open problem}.

\begin{remark}[Stochastic Lanczos quadrature]
One promising direction is \emph{stochastic Lanczos quadrature} (SLQ):
estimate $\log|K|$ as $m\cdot\E_{\mathbf{u}}[\log\tilde\mu(\mathbf{u})]$,
where $\tilde\mu$ is a stochastic Gauss quadrature approximation to the
spectral measure of $K$, requiring only matrix--vector products with $K$
(which remain $O(m)$).  The gradient \eqref{eq:grad-logdet-aniso} can be
obtained by differentiating the Gauss nodes and weights with respect to $H$.
This approach is left for future work.
\end{remark}

\bibliographystyle{plainnat}
\bibliography{refs}

\appendix

\section{Lessons Learned and Paths Avoided}
\label{sec:lessons}

This appendix records design decisions that were made, approaches that were
tried and abandoned, and reasons why certain seemingly natural paths turned out
to be dead ends.  It is intended as a living lab-notebook entry: when a
promising direction is ruled out, the reasoning should be recorded here so it
need not be rediscovered later.

\subsection{Delaunay Triangulation Does Not Scale to Higher Dimensions}
\label{sec:lesson-delaunay}

The Level-1 method (Section~\ref{sec:baseline}) builds a Delaunay
triangulation of the domain and assembles FEM matrices over the resulting
simplices.  Two obstacles make this approach untenable in dimensions $D \ge 3$
or even in $D=2$ once adaptive resolution is required.

\paragraph{Combinatorial explosion of simplices.}
The standard Freudenthal--Kuhn triangulation of $[0,1]^D$ decomposes the
hypercube into exactly $D!$ simplices, one per permutation of the coordinate
axes~\citep{freudenthal1942}.  More efficient triangulations exist: the best
known construction achieves roughly $(0.816)^D \cdot D!$ simplices
\citep{ordensantos2003}, and current lower bounds~\citep{glazyrin2012} are
weaker still.  The precise minimum (the \emph{simplexity} of the $D$-cube)
remains an open problem in combinatorial geometry.  In any case all known
constructions and lower bounds grow super-exponentially in $D$: for $D=4$ the
Freudenthal triangulation already uses $4! = 24$ top-dimensional simplices per
unit hypercube, so a grid of $n^D$ points produces $O(n^D \cdot D!)$ simplices
and FEM matrix assembly over all of them is prohibitive.

\paragraph{Non-stationarity of the triangulation under refinement.}
Adaptive resolution requires adding new vertices to the mesh.  Inserting a
point into a Delaunay triangulation triggers a Bowyer--Watson cascade of
simplex flips whose extent is unpredictable (it depends on the local geometry).
This means the mesh is \emph{not static}: the full triangulation, and hence the
FEM matrices $C$ and $G$, must be recomputed from scratch after every RJMCMC
birth or death move.  The cost per move is therefore $O(m \log m)$ just for
retriangulation, before any matrix assembly.

\paragraph{Conclusion.}
Delaunay triangulation is a natural first choice but hits a fundamental
combinatorial wall in higher dimensions and does not compose well with adaptive
sampling.  We abandon it in favour of the hat-function tree basis described in
Section~\ref{sec:hat-tree}.

\subsection{Hat-Function Tree Basis as the Natural Replacement}
\label{sec:lesson-tree-basis}

The adaptive hierarchical basis (Section~\ref{sec:hat-tree}) resolves both
obstacles above.

\begin{itemize}
  \item \textbf{No combinatorial explosion.}
    A $D$-dimensional basis function is the product of $D$ one-dimensional
    functions (constant, linear, or hat).  Adding a node to the tree requires
    only $O(D)$ work to compute its states and at most $O(D)$ new inner
    products to update $C$ and $G$; no global retriangulation is needed.
  \item \textbf{Static structure under refinement.}
    The mass and stiffness matrices have a \emph{recursive block structure}
    tied to the tree (Section~\ref{sec:hat-mass}).  A birth move in dimension
    $d$ splits one block into sub-blocks; existing off-diagonal blocks are
    unchanged (sibling hats have zero mass overlap, Eq.~\eqref{eq:C-block-zero}).
    The stiffness matrix $G$ is diagonal and stays diagonal after any birth
    (Section~\ref{sec:hat-stiffness}).  FEM matrix updates are therefore local
    and $O(1)$ per move.
  \item \textbf{Natural low-to-high resolution ordering.}
    The tree encodes a multi-resolution hierarchy: the root is the constant
    function~1, and successive refinements add finer and finer hat functions.
    This gives a principled prior over resolution and a natural initialisation
    (start at the root, i.e., a constant intensity field) rather than requiring
    an expensive high-resolution mesh as the initial condition.
\end{itemize}


\end{document}
